% ============================================================
% Paper 3 — A Finite-Cutoff Hilbert-Space Model
%           for Prime--Zero Energy Structure
%
% Author: Ulrich Tehrani
% Version: v3.2 (May~2026)
%
% Changelog
% ---------
% v1 (March 2026)
%   Initial publication.
%
% v2 (April~2026)
%   Canonical weights updated to normative form
%   c_p = sqrt(V_p(1/2)). Asymptotic H_local(1/2,kappa)
%   corrected from 8(log kappa)^2 to 2(log kappa)^2.
%   All formal proofs and numerical results unchanged.
%
% v3 (April~2026)
%   Series consolidation. Template v1.7 applied:
%   Notation section added, bibliography scoped and
%   harmonized, footer format updated, source hygiene.
%   Structural cleanup: Non-Circularity moved to §1.2,
%   Call for Feedback as standalone section, Variable
%   Reference replaced by Normative Parameter Table,
%   §1.5 expanded with positioning (Berry-Keating,
%   de Branges, Goldfeld cited). Hardcoded cross-paper
%   theorem numbers replaced by section references.
%   No mathematical content changed.
%
% v3.2 (May~2026)
%   Canonical weight notation corrected (audit, May 2026):
%   c_p = sqrt(V_p(1/2)) was a verbal mismatch. The exact
%   formula 4(log p)^2(2p-1)/(p(p-1)^2) is f_p, not V_p(1/2).
%   The canonical weight in the eta framework is now written
%   c_p^eta := sqrt(f_p), numerically equal to Paper 2's
%   c_p^ren = sqrt(f_p) but with a distinct role inside the
%   eta_orig formula. All numerical values unchanged.
% ============================================================

\documentclass[11pt]{article}
\usepackage{cmap}            % ToUnicode CMaps for §, en/em-dashes etc.
\usepackage[T1]{fontenc}
\emergencystretch=2em        % avoid overfull hboxes from inline math
\usepackage[utf8]{inputenc}
\usepackage[a4paper,margin=1in]{geometry}
\usepackage{amsmath,amssymb,amsthm,mathtools}
\usepackage{xcolor}
\usepackage{hyperref}
\usepackage{enumitem}
\usepackage{setspace}
\usepackage{booktabs}
\usepackage{array}
\usepackage{graphicx}

\hypersetup{colorlinks=true,linkcolor=blue!50!black,urlcolor=blue!50!black,citecolor=blue!50!black}
\setlength{\parskip}{0.5em}
\setlength{\parindent}{0pt}
\onehalfspacing

% -- Theorem environments 
\newtheorem{theorem}{Theorem}[section]
\newtheorem{proposition}[theorem]{Proposition}
\newtheorem{lemma}[theorem]{Lemma}
\newtheorem{corollary}[theorem]{Corollary}
\newtheorem{definition}[theorem]{Definition}
\newtheorem{assumption}[theorem]{Assumption}
\newtheorem{observation}[theorem]{Observation}
\newtheorem*{remark}{Remark}
\newtheorem{openproblem}[theorem]{Open Problem}

% -- Macros 
\newcommand{\Hstr}{H_{\mathrm{str}}}
\newcommand{\Hnull}{H_{\mathrm{null}}}
\newcommand{\Estr}{E_{\mathrm{str}}}
\newcommand{\Espec}{E_{\mathrm{spec}}}
\newcommand{\etaorig}{\eta_{\mathrm{orig}}}
\newcommand{\Tren}{T_{\mathrm{ren}}}
\newcommand{\Gun}{G^{\mathrm{un}}}
\newcommand{\Gnorm}{G^{\mathrm{norm}}}
\newcommand{\acan}{a_p^{\mathrm{can}}}
\newcommand{\Dcan}{D^{\mathrm{can}}}
\newcommand{\Estrcan}{\Estr^{\mathrm{can}}}
\newcommand{\Especcan}{\Espec^{\mathrm{can}}}
\newcommand{\Hlocal}{H_{\mathrm{local}}}
\newcommand{\Hren}{H_{\mathrm{ren}}}
\newcommand{\cren}{c_p^{\mathrm{ren}}}
\newcommand{\ceta}{c_p^{\eta}}
\newcommand{\fp}{f_p}

% -- Title 
\title{\vspace{-1.5em}\bfseries A Finite-Cutoff Hilbert-Space Model\\[0.4em]
\large for Prime--Zero Energy Structure}
\author{Ulrich Tehrani}
% FIX v3.2 — c_p^eta canonical form (HIGH 4 audit, May 2026)
\date{Research Note \textperiodcentered\ May~2026 \textperiodcentered\ v3.2}

\begin{document}
\maketitle

\begin{abstract}
We construct a finite-cutoff Hilbert-space model for the spectral structure
associated with primes and nontrivial zeros of the Riemann zeta function.
Given a prime cutoff $\kappa$ and truncated zero ordinates
$\{\gamma_1,\ldots,\gamma_N\}$, we define two finite-dimensional real
Hilbert spaces $\Hstr$ (over primes) and $\Hnull$ (over zeros),
connected by a linear map $\Phi\colon\Hstr\to\Hnull$, and the
self-adjoint loop operator $T = \Phi^*\circ\Phi$.

\textbf{What is proved.}
We establish the algebraic identity
$\Delta := \Estr - \Espec
 = \sum_p c_p^2\|a_p\|^2 - \|\sum_p c_p a_p\|^2$
by direct expansion.

\textbf{What is numerical.}
$\Delta > 0$ and
$\etaorig = \Delta/\Estr \in [0.598,\,0.704]$ for benchmark parameters
$(\kappa,\varepsilon,N)=(53,0.05,100)$, $\sigma\in[0.1,0.95]$.

\textbf{What remains open.}
An analytic proof of $\etaorig>0$ remains open.

\textbf{Structure.}
The formal core in \S\S\,2--5 is non-circular and proof-complete;
results from Papers~1--2 appear only as context.
No use of the Riemann Hypothesis, the GUE conjecture,
or the Hilbert--P\'olya postulate is made anywhere in \S\S\,2--5.
The model suggests a picture in which $\sigma=\tfrac12$ plays
the role of an energy reference point.
\end{abstract}

\vspace{2em}

% --------------------------------------------------------
% Notation and Conventions  (template invariant, v1.3)
% --------------------------------------------------------
\section*{Notation and Conventions}
\label{sec:notation}

\emph{Critical line.}
Throughout this series, $\sigma$ denotes $\Re(s)$ and the critical
line means $\sigma=\tfrac12$. The functional equation symmetry
$\xi(s)=\xi(1-s)$ makes $\sigma=\tfrac12$ the natural reference
point.

\medskip
\emph{Global parameters.}
In this paper, $\kappa=53$ is the prime cutoff (16~active primes
$p\leq\kappa$), $\varepsilon=0.05$ the Gaussian regularization
parameter, and $N=100$ the number of Riemann zero ordinates
$0<\gamma_1\leq\gamma_2\leq\cdots\leq\gamma_N$ employed.

\medskip
\emph{Main objects.}
The prime space $\Hstr=\ell^2(S(\kappa))$ and zero space
$\Hnull=\ell^2(\{\gamma_1,\ldots,\gamma_N\})$ are connected by
$\Phi\colon\Hstr\to\Hnull$, $\Phi(e_p)=a_p$, where
$(a_p)_k=e^{-\varepsilon^2\gamma_k^2/2}\sin(\gamma_k\log p)$.
The loop operator is $T=\Phi^*\circ\Phi$ and
$\etaorig=1-\Espec/\Estr$ measures the energy gap.

\medskip
\emph{Weights.}
The renormalized weight squared
\[
  \fp \;:=\; V_p(\tfrac12) - \frac{4(\log p)^2}{p}
       \;=\; \frac{4(\log p)^2(2p-1)}{p\,(p-1)^2}
\]
subtracts the leading divergent term from the local prime curvature
$V_p(\tfrac12)$ of Paper~1. The canonical weight in the
$\eta$-framework of this paper is
\[
  \ceta \;:=\; \sqrt{\fp}\,.
\]

\medskip
\emph{Not to be confused with.}
$T=\Phi^*\Phi$ (loop operator on $\Hstr$, this paper) is distinct from
$\tilde T=\Phi\Phi^*$ (the adjoint loop operator on $\Hnull$).
$D\approx 9.471$ (renormalized energy sum, Paper~2) is distinct from
$\Estr\approx 6.119$ (structural energy at benchmark parameters,
this paper). The canonical weight $\ceta = \sqrt{\fp}$ in the
$\eta$-framework of this paper has the same numerical value as
Paper~2's renormalized weight $\cren = \sqrt{\fp}$ but plays a
distinct role: $\ceta$ enters $\Estr=\sum_p (\ceta)^{2}\|a_p\|^2$
inside the $\eta$-formula, whereas $\cren$ enters
$D=\sum_p (\cren)^{2}$ in the Weil renormalization.

\newpage

%% ═══════════════════════════════════════════════════════════════════════════
\section{Introduction}
%% ═══════════════════════════════════════════════════════════════════════════

\subsection{Background and Motivation}

The present paper is the third in a series developing an
operator-theoretic framework for studying the distribution of the
non-trivial zeros of the Riemann zeta function $\zeta(s)$.
The formal core in \S\S\,2--4 is non-circular and proof-complete.

\textbf{Paper~1}~\cite{Paper1} introduced the curvature function
$H_{\xi}(\sigma,t) := \frac{d^2}{d\sigma^2}\log|\xi(\sigma+it)|^2$
and established the divergence
$\Hlocal(\tfrac12,\kappa)\sim 2(\log\kappa)^2\to\infty$ at
$\sigma=\tfrac12$, contrasted with convergence for $\sigma>\tfrac12$.

\textbf{Paper~2}~\cite{Paper2} connected this local curvature to the
Weil explicit functional, constructing admissible test functions
and establishing the convergent weight sum
$D=\sum_p(c_p^{\mathrm{ren}})^2\approx 9.471$.
These results appear in \S\,2 for motivation and context;
the formal core does not depend on them.

\textbf{The present paper} makes the geometric structure underlying
both papers explicit: two finite-dimensional Hilbert spaces $\Hstr$
and $\Hnull$, connected by a linear map $\Phi$, and a self-adjoint
loop operator $T=\Phi^*\circ\Phi$.

Several established programs seek a spectral interpretation of the
Riemann zeros. Berry and Keating~\cite{BerryKeating} studied a
semiclassical analogue via quantum chaos;
Connes~\cite{Connes} realizes the zeros through a trace formula in
noncommutative geometry;
de~Branges~\cite{deBranges} approached the problem through Hilbert
spaces of entire functions. The present model differs from all three:
it constructs \emph{two} finite-dimensional spaces coupled by a
bilinear map, rather than seeking a single operator whose spectrum is
$\{\gamma_k\}$. The Gram-matrix construction
$T_{pq}=\langle a_p,a_q\rangle$ from damped sine vectors appears to
be new; conceptually it may be viewed as a discrete analogue of
spectral interpretations of the explicit formula, as studied by
Goldfeld~\cite{Goldfeld}. The test functions of Paper~2 connect to
adelic Schwartz functions as in Tate's thesis~\cite{Tate}; the Weil
positivity criterion~\cite{Weil,Lagarias} provides the outer
framework.

The renormalized framework of Paper~2 and the $\eta$-framework of
this paper use different normalizations; a precise identification
remains open (Open Problem~\ref{op:shell}).

\subsection{Main Results}

We define two finite-dimensional real Hilbert spaces $\Hstr$ (over
primes $p\le\kappa$) and $\Hnull$ (over zero ordinates
$\gamma_1,\ldots,\gamma_N$), connected by the linear map
$\Phi:\Hstr\to\Hnull$ built from Gaussian-damped prime resonance
vectors (\S\,3). The self-adjoint loop operator $T=\Phi^*\circ\Phi$
is the Gram matrix of these vectors (\S\,3).

We prove the algebraic identity $\Delta:=\Estr-\Espec$ by direct
expansion (Theorem~4.3). Numerically, $\etaorig=\Delta/\Estr\in
[0.598,0.704]$ for benchmark parameters
$(\kappa,\varepsilon,N)=(53,0.05,100)$, $\sigma\in[0.1,0.95]$
(Observation~\ref{obs:nonneg}). An analytic proof of $\etaorig>0$ remains open
(Open Problem~\ref{op:pos}).

\subsection{Reader Contract}

This paper \emph{proves}, by direct algebraic expansion, the
variance identity $\Delta=\Estr-\Espec$ (Theorem~4.3) and the basic
properties of the loop operator $T=\Phi^*\circ\Phi$: self-adjointness,
positive semi-definiteness, and $\sigma$-invariance of the normalized
Gram matrix (Proposition~3.5). These are unconditional.

This paper \emph{establishes numerically} that $\etaorig>0$ for all
tested parameter combinations, with
$\etaorig\in[0.598,0.704]$ at benchmark parameters (\S\,5).
The uniform spectral bound $B_{\max}\le 0.106<1$ and its
$\sigma$-invariance are numerical observations.

This paper \emph{leaves open} the analytic proof of $\etaorig>0$
(Problem~6.1), the spectral bound for canonical weights
(Problem~6.2), the analytical cancellation bound (Problem~6.3),
the $\sigma$-invariance of $B_{\max}$ (Problem~6.4), and the
renormalized shell energy (Problem~6.5).

%% ═══════════════════════════════════════════════════════════════════════════
\section{Two Hilbert Spaces}
%% ═══════════════════════════════════════════════════════════════════════════

We work throughout with real Hilbert spaces; all vectors and operators have
real entries. Complex phasor language appears in \S\,3.1 as
heuristic motivation only; all formal objects ($\Hstr, \Hnull, \Phi, T$)
are real.

Throughout this paper, $\kappa>1$ is a fixed real cutoff parameter,
$S(\kappa):=\{p\text{ prime}:p\leq\kappa\}$ the set of active primes,
and $\{\gamma_k\}_{k=1}^N$ the first $N$ positive imaginary parts of
nontrivial zeros of $\zeta(s)$, ordered $0<\gamma_1\leq\gamma_2\leq\cdots$.
The Gaussian damping parameter $\varepsilon>0$ is fixed throughout.

\begin{definition}[Prime space $\Hstr$]
\[
\Hstr := \ell^2(S(\kappa))
\]
with basis $\{e_p : p\in S(\kappa)\}$, standard inner product
$\langle a,b\rangle_{\mathrm{str}}=\sum_p a_p b_p$, and
$\dim\Hstr=|S(\kappa)|=\pi(\kappa)$.
The \emph{structural energy} of a weight vector $c\in\Hstr$ is
$\Estr(c)=\sum_p c_p^2\|a_p\|^2$ (defined in \S4 after $\Phi$ is introduced).
\end{definition}

\begin{definition}[Zero space $\Hnull$]
\[
\Hnull := \ell^2(\{\gamma_1,\ldots,\gamma_N\})
\]
with basis $\{f_k:k=1,\ldots,N\}$, inner product
$\langle u,v\rangle_{\mathrm{null}}=\sum_{k=1}^N u_k v_k$
(unweighted), and $\dim\Hnull=N$.
\end{definition}

\begin{remark}[Why $\sigma=\tfrac12$ is a distinguished reference point]
\label{rem:half}
This remark provides context; it is not used in \S\S\,3--5.

Two independent proven reasons support this choice:

\emph{(i) Symmetry} [proven]: $\sigma=\tfrac12$ is the unique fixed point of
$s\mapsto 1-s$, the symmetry of the completed zeta function
$\xi(s)=\xi(1-s)$.

\emph{(ii) Curvature singularity} [proven, Papers~1--2]:
$\Hlocal(\tfrac12,\kappa)\sim 2(\log\kappa)^2\to\infty$ as $\kappa\to\infty$,
while $\Hlocal(\sigma,\kappa)\to C(\sigma)<\infty$ for all $\sigma>\tfrac12$.

$\sigma=\tfrac12$ is thus a distinguished reference point by symmetry and
curvature divergence. Energy near-equality ($\Estr\approx\Espec$) is NOT
observed in the finite-cutoff model (see \S5 and \S8).
\end{remark}

%% ═══════════════════════════════════════════════════════════════════════════
\section{The Linear Map $\Phi$ and the Loop Operator $T$}
%% ═══════════════════════════════════════════════════════════════════════════

\emph{Convention.} Gaussian damping is carried by the vectors:
$a_p=(e^{-\varepsilon^2\gamma_k^2/2}\sin(\gamma_k\log p))_k$.
The $\Hnull$ inner product is unweighted: $\langle u,v\rangle=\sum_k u_kv_k$.

\subsection{Motivation for the resonance vectors}

The formal definition is Definition~3.2 below; this subsection
provides intuition.

The prime resonance vector at zero ordinate $\gamma_k$ and prime $p$ arises as
the imaginary part of the prime phasor $\varphi_p(u)=e^{iu\log p}$ evaluated
at $u=\gamma_k$, with Gaussian regularization for convergence:
\[
(a_p)_k = e^{-\varepsilon^2\gamma_k^2/2}\cdot\sin(\gamma_k\log p).
\]
The $k$-th component measures how strongly prime $p$ resonates with zero
ordinate $\gamma_k$.

\begin{definition}[Prime resonance vectors]
For each prime $p\in S(\kappa)$, the \emph{canonical resonance vector}
at $\sigma=\tfrac12$ is
\[
\acan := a_p :=
\bigl(e^{-\varepsilon^2\gamma_k^2/2}\sin(\gamma_k\log p)\bigr)_{k=1}^N
\;\in\;\Hnull.
\]
At general $\sigma$, the resonance vectors depend on $\sigma$ via
\[
a_p^{(\sigma)} := p^{-(\sigma-1/2)}\cdot\acan.
\]
The vectors $a_p=\acan=a_p^{(1/2)}$ at $\sigma=\tfrac12$ are the
canonical representatives used in all primary definitions.
\end{definition}

\begin{definition}[Linear map $\Phi$]
\[
\Phi\colon\Hstr\to\Hnull,\qquad\Phi(e_p):=\acan.
\]
Extended linearly: $\Phi(c)=\sum_p c_p\acan$ for any
$c=\sum_p c_p e_p\in\Hstr$.

\emph{We define $\Phi$ as a linear map; no injectivity claim is made for the
finite-cutoff model.}
\end{definition}

\begin{definition}[Gram matrices]
The \emph{unnormalized Gram matrix} (canonical) is
\[
\Gun_{pq} := \Gun{}^{\mathrm{,can}}_{pq}
:= \langle\acan,a_q^{\mathrm{can}}\rangle_{\Hnull}
= \sum_{k=1}^N e^{-\varepsilon^2\gamma_k^2}\sin(\gamma_k\log p)\sin(\gamma_k\log q).
\]
The $\sigma$-dependent extension is
$\Gun(\sigma)_{pq}=p^{-(\sigma-1/2)}q^{-(\sigma-1/2)}\Gun_{pq}$.
The \emph{normalized Gram matrix} is
\[
\Gnorm_{pq}:=\frac{\Gun_{pq}}{\|\acan\|\cdot\|a_q^{\mathrm{can}}\|}
=\Bigl\langle\frac{\acan}{\|\acan\|},\frac{a_q^{\mathrm{can}}}{\|a_q^{\mathrm{can}}\|}\Bigr\rangle_{\Hnull}.
\]
The matrix $\Gnorm$ is $\sigma$-invariant (Proposition~3.5(iv)).
\end{definition}

\begin{definition}[Loop operator $T$]
\[
T:=\Phi^*\circ\Phi\colon\Hstr\longrightarrow\Hstr.
\]
The matrix representation of $T$ in the prime basis is $T_{pq}=\Gun_{pq}$.
$T$ is thus identical to the canonical Gram matrix $\Gun$, and is explicitly
constructed from prime resonance vectors.
\end{definition}

\begin{proposition}[Properties of $T$; proved]
\label{prop:T}
\begin{enumerate}[label=\emph{(\roman*)}, leftmargin=2.2em]
\item \emph{Self-adjoint:} $T^*=(\Phi^*\Phi)^*=\Phi^*\Phi=T$.
\item \emph{Positive semi-definite:} $\langle Tc,c\rangle=\|\Phi c\|^2\geq 0$
for all $c\in\Hstr$.
\item \emph{Non-negative eigenvalues:} All eigenvalues $\lambda_j\geq 0$.
\item \emph{$\sigma$-invariance of $\Gnorm$} [by direct calculation]:
For the $\sigma$-dependent family $a_p^{(\sigma)}$, the normalized Gram matrix
$\Gnorm_{pq}$ is independent of $\sigma$. \emph{Proof.} The scalar factor
$p^{-(\sigma-1/2)}$ cancels in the normalized ratio. \qed
\end{enumerate}
\end{proposition}

\begin{remark}[Spectral structure of $T$]
The eigenvalues $\lambda_j\geq 0$ of $T=\Phi^*\circ\Phi$ are the squares of
the singular values $\sigma_j(\Phi)$: $\lambda_j=\sigma_j(\Phi)^2$.
The eigenvalues $\lambda_j$ are built from the zero ordinates $\gamma_k$
appearing in the definition of $\acan$, but are \emph{not} directly the zeros
$\gamma_k$ themselves. This distinguishes $T$ from the classical Hilbert-P\'olya
ansatz: $T$ is explicitly constructed, not postulated.
\end{remark}

%% ═══════════════════════════════════════════════════════════════════════════
\section{The Variance Identity and $\etaorig$}
%% ═══════════════════════════════════════════════════════════════════════════

\emph{Convention.} Objects without $\sigma$-argument refer to the canonical
($\sigma=\tfrac12$) version throughout \S\,4.

\begin{definition}[Canonical and $\sigma$-dependent energies]
\label{def:energies}
For a weight vector $c\in\Hstr$ with components $c_p$:

\emph{Canonical quantities} (at $\sigma=\tfrac12$):
\[
\Dcan:=\mathrm{diag}\!\bigl(\|\acan\|^2\bigr)_{p\in S(\kappa)},
\qquad
\Estrcan:=c^\top\Dcan c=\sum_p c_p^2\|\acan\|^2,
\]
\[
\Especcan:=\|\Phi c\|^2=c^\top\Gun c.
\]

\emph{$\sigma$-dependent extensions:}
\[
D(\sigma)_{pp}:=p^{-2(\sigma-1/2)}\cdot\Dcan_{pp},
\qquad
\Estr(\sigma):=\sum_p c_p^2\,p^{-2(\sigma-1/2)}\|\acan\|^2,
\]
\[
\Espec(\sigma):=c^\top\Gun(\sigma)\,c.
\]

The \emph{canonical weight vector} in the $\eta$-context is
\[
\ceta \;:=\; \sqrt{\fp}
    \;=\; \sqrt{\frac{4(\log p)^2(2p-1)}{p(p-1)^2}},
    \qquad p \in S(\kappa).
\]
\begin{remark}[Canonical weights --- normative form]
The canonical weight $\ceta = \sqrt{\fp}$
is the normative form used throughout this series. For large $p$,
$\ceta \approx \sqrt{8}\cdot(\log p)/p$; the two forms
agree asymptotically but differ for small primes.
The numerical results in \S\,5 use the normative form.
\end{remark}
\end{definition}

\begin{assumption}[Positive definiteness of $D$]
Throughout this paper we assume $D_{pp}=\|a_p\|^2>0$ for all
$p\in S(\kappa)$. This holds in all numerical experiments conducted here.
\end{assumption}

\begin{theorem}[Algebraic identity; proved]
\label{thm:identity}
The following equality holds by direct expansion:
\[
\Delta:=\Estr-\Espec=\sum_p c_p^2\|a_p\|^2-\Bigl\|\sum_p c_p a_p\Bigr\|^2.
\]
\end{theorem}

\begin{proof}
Expand $\Espec=\|\Phi c\|^2=\bigl\langle\sum_p c_p a_p,\sum_q c_q
a_q\bigr\rangle=\sum_{p,q}c_p c_q\langle a_p,a_q\rangle=c^\top\Gun c$.
Then $\Delta=c^\top\Dcan c-c^\top\Gun c$, which is the stated identity.
\end{proof}

\begin{observation}[Non-negativity, numerical]
\label{obs:nonneg}
\[
\Delta\geq 0\quad\text{equivalently}\quad\etaorig\geq 0
\]
holds for all tested parameters with canonical weights
$\ceta=\sqrt{\fp}$ (Definition~\ref{def:energies}).
\emph{[Numerical only --- no unconditional proof for general $c$.]}

Note: $\etaorig\geq 0$ for general weight vector $c$ is equivalent to
$\lambda_{\max}(D^{-1/2}TD^{-1/2})\leq 1$, which is the content of
Open Problem~\ref{op:spectral}. Numerical experiments show
$\lambda_{\max}(D^{-1/2}TD^{-1/2})\approx 0.39\cdot\pi(\kappa)\to\infty$,
so the universal bound $\lambda_{\max}\leq 1$ does not hold in general.

\emph{Cauchy-Schwarz does \textbf{not} guarantee $\etaorig\geq 0$}: the
correct identity is $\Delta=c^\top(\Dcan-\Gun)c$; the sign of $\Delta$ is not
determined by the algebraic identity alone.
\end{observation}

\begin{definition}[$\etaorig$]
\label{def:eta}
\[
\etaorig:=1-\frac{\Espec}{\Estr}=1-\frac{\|\Phi c\|^2}{\Estr}.
\]
Boundary cases: $\etaorig=0$ iff $\Espec=\Estr$; $\etaorig=1$ iff $\Phi c=0$.
\end{definition}

\begin{remark}[$\sigma$-dependence via rescaled weights]
The entire $\sigma$-dependence of $\etaorig(\sigma)$ is carried by the
rescaled weight vector $\tilde{c}_p(\sigma):=c_p\cdot p^{-(\sigma-1/2)}$:
\[
\etaorig(\sigma)=1-\frac{\tilde{c}(\sigma)^\top\Gun\,\tilde{c}(\sigma)}
{\tilde{c}(\sigma)^\top\Dcan\,\tilde{c}(\sigma)}.
\]
\emph{Proof.} Since $\Gun(\sigma)_{pq}=p^{-(\sigma-1/2)}q^{-(\sigma-1/2)}\Gun_{pq}$,
one has $c^\top\Gun(\sigma)c=\tilde{c}^\top\Gun\tilde{c}$ and
$\Estr(\sigma)=\tilde{c}^\top\Dcan\tilde{c}$. \qed
\end{remark}

\emph{Numerical evidence.}
All computations use $N=100$ zero ordinates at the reference
parameters; the algebraic identity (Theorem~\ref{thm:identity}) is
verified to rounding error $<10^{-14}$. Verification scripts are
available in the companion repository~\cite{Code}.

\begin{observation}[$\etaorig$ positive across all tested parameters]
$\etaorig>0$ in all tested parameter combinations. For the benchmark slice
$(\kappa,\varepsilon,N)=(53,0.05,100)$, $\sigma\in[0.1,0.95]$:
\[
\etaorig\in[0.598,\;0.704].
\]
The function $\sigma\mapsto\etaorig(\sigma)$ is strictly monotone decreasing.
Note: the range $[0.598,0.704]$ is strongly $\varepsilon$-dependent; no
universal range over all $(\kappa,\varepsilon)$ is claimed.
\end{observation}

The following values illustrate the strict monotone decrease of
$\etaorig(\sigma)$ at the reference parameters.

\begin{center}
\begin{tabular}{cc}
\toprule
$\sigma$ & $\etaorig$ \\
\midrule
0.10 & 0.70441977 \\
0.20 & 0.69826865 \\
0.30 & 0.69038398 \\
0.40 & 0.68072993 \\
0.50 & 0.66926873 \\
0.60 & 0.65602069 \\
0.70 & 0.64110102 \\
0.80 & 0.62472010 \\
0.90 & 0.60715692 \\
0.95 & 0.59802896 \\
\bottomrule
\end{tabular}
\end{center}

\begin{observation}[Uniform spectral bound $B_{\max}\leq 0.106<1$]
Define the mode contribution
$B_j(\sigma):=\lambda_j(\sigma)|\langle c,u_j(\sigma)\rangle|^2/\Estr(\sigma)$
with consistency identity $\sum_j B_j(\sigma)=1-\etaorig(\sigma)$ [algebraic].
The empirical diagnostic $B_{\max}(\sigma):=\max_{j:\lambda_j>1}B_j(\sigma)$
satisfies
\[
B_{\max}(\kappa,\varepsilon,\sigma)\leq 0.106<1
\]
for all tested $(\kappa,\varepsilon,\sigma)$.
\end{observation}

Selected values:

\begin{center}
\begin{tabular}{cccc}
\toprule
$\kappa$ & $\varepsilon=0.02$ & $\varepsilon=0.05$ & $\varepsilon=0.10$ \\
\midrule
23  & 0.1047 & 0.0220 & 0.0000 \\
53  & 0.0793 & 0.0737 & 0.0000 \\
101 & 0.0537 & 0.0844 & 0.0324 \\
199 & 0.0891 & 0.0461 & 0.1058 \\
\bottomrule
\end{tabular}
\end{center}

\begin{observation}[$B_{\max}(\sigma)$ numerically $\sigma$-invariant]
For $(\kappa,\varepsilon)=(53,0.05)$, $B_{\max}(\sigma)=0.0737$ for all
tested $\sigma\in\{0.3,0.4,0.5,0.6,0.7\}$. Analytical explanation: open
(Open Problem~\ref{op:sigma}).
\end{observation}

\begin{remark}[Cancellation mechanism, numerical analysis]
The mechanism underlying $B_{\max}<1$ is a cancellation between the monotone

weight vector $\ceta=\sqrt{\fp}$ and the oscillatory eigenvectors $u_j$
of $\Gun$: for $\kappa=53$, eigenvectors $u_1,\ldots,u_4$ each exhibit 10
sign changes along the prime ordering. Cancellation rates 67--97\% observed
for $\kappa\in\{23,53,101,199\}$. Formal proof: Open Problem~\ref{op:cancel}.
\end{remark}

%% ═══════════════════════════════════════════════════════════════════════════
\section{Open Problems}
%% ═══════════════════════════════════════════════════════════════════════════


\begin{openproblem}[Analytic positivity]\label{op:pos}
Prove analytically that $\etaorig>0$ for the canonical weight vector
$\ceta=\sqrt{\fp}$.
Numerically confirmed: $\etaorig\in[0.598,0.704]$ for $(\kappa,\varepsilon,N)=(53,0.05,100)$,
$\sigma\in[0.1,0.95]$ (\S\,4).
\end{openproblem}

\begin{openproblem}[Spectral bound for canonical weights]\label{op:spectral}
Let $\Tren:=D^{-1/2}TD^{-1/2}$.
Numerical experiments show $\lambda_{\max}(\Tren)\approx 0.39\cdot\pi(\kappa)\to\infty$;
the universal bound $\lambda_{\max}<1$ (for all $c$) does not hold.
However, for the canonical weight vector $\ceta=\sqrt{\fp}$,
the quantity $\langle\Tren\tilde{c},\tilde{c}\rangle$ satisfies
$\langle\Tren\tilde{c},\tilde{c}\rangle<1$ (equivalently $\etaorig>0$)
for all tested $\kappa\in\{23,53,101,199,503,1009\}$.
Prove this analytically. The canonical $c_p$ appears orthogonal to the
dominant eigenmode of $\Tren$ (localized at $\gamma_1=14.13\ldots$);
proving this orthogonality is the core open problem.
\end{openproblem}

\begin{openproblem}[Analytical cancellation bound]\label{op:cancel}
Prove analytically:
\[
\max_{j:\,\lambda_j>1}\frac{|\langle c,u_j\rangle|^2}{\|c\|^2}\leq C_1\ll 1
\quad\text{uniformly in }\kappa,j.
\]
Approach: Abel summation for $c_p\sim(\log p)/p$ against oscillatory partial
sums $\sum_{q\leq p}(u_j)_q$, without reference to RH.
Numerically: $C_1<0.36$ for $(\kappa,\varepsilon)=(53,0.05)$.
\end{openproblem}

\begin{openproblem}[$\sigma$-invariance of $B_{\max}$ analytically]\label{op:sigma}
The numerical $\sigma$-invariance of $B_{\max}(\sigma)$ (consistent with
Proposition~\ref{prop:T}(iv)) is not yet derived analytically from the weight
rescaling $c\mapsto\tilde{c}(\sigma)$.
\end{openproblem}


\begin{openproblem}[Renormalized shell energy]\label{op:shell}
The renormalized shell energy $\Hren(\kappa,\varepsilon(\kappa))$ for canonical
coupling $N_{\mathrm{eff}}=\pi(\kappa)$ is not yet computed. This lies in the
Paper~2 normalization framework.
\end{openproblem}

%% ═══════════════════════════════════════════════════════════════════════════
\section{Non-Circularity}
\label{sec:nocirc}
%% ═══════════════════════════════════════════════════════════════════════════

We document explicitly that the results of this paper do not rely on
the objects they are directed toward.

\textbf{No RH as input.}
At no point is the Riemann Hypothesis assumed.
The zero ordinates $\gamma_k$ enter the construction of $\Phi$ as
numerically specified inputs, and no property of their distribution
requiring the Riemann Hypothesis is assumed or used.

\textbf{No GUE, Montgomery, or Hilbert--P\'olya hypothesis.}
The loop operator $T=\Phi^*\Phi$ is built explicitly from the prime
resonance vectors $\{a_p\}_{p\le\kappa}$. Its construction postulates
neither GUE statistics, Montgomery's pair correlation conjecture, nor
the Hilbert--P\'olya conjecture. The eigenvalues $\lambda_j$ of $T$ are
the squares of the singular values of $\Phi$ and are not identified
with zero ordinates $\gamma_k$.

\textbf{Status separation maintained.}
The variance identity $\Delta=\Estr-\Espec$ is an algebraic theorem;
the sign statement $\Delta\geq 0$ is established numerically on the
tested grid. These are distinct levels and are flagged separately
throughout. No claim relating $\etaorig$ to RH is proved or used.

%% ═══════════════════════════════════════════════════════════════════════════
\section*{Acknowledgments}
%% ═══════════════════════════════════════════════════════════════════════════

The author thanks the anonymous readers who provided feedback on
earlier versions of this paper.

%% ═══════════════════════════════════════════════════════════════════════════
\section*{Call for Feedback}

This paper is part of an ongoing open research initiative toward the
Riemann Hypothesis.
All papers in the series are freely available on Zenodo under CC~BY~4.0,
and the accompanying verification code is hosted on GitHub.

{\sloppy
Zenodo (\url{https://zenodo.org/search?q=Ulrich+Tehrani}) and GitHub
(\url{https://github.com/utehrani/analysislab-nt}).
\par}

Topics of particular interest include: whether $\etaorig>0$ is an
expected consequence of known equidistribution results for the phases
$\{\gamma_k\log p \bmod 2\pi\}$, and whether the observed
near-orthogonality of the weight vector to the dominant eigenmode of
$T$ connects to standard cancellation phenomena in partial sums
involving primes and zeros.

%% ═══════════════════════════════════════════════════════════════════════════

\newpage
\begin{thebibliography}{99}

\bibitem{BerryKeating}
M.V.~Berry and J.P.~Keating,
``The Riemann zeros and eigenvalue asymptotics'',
\emph{SIAM Rev.} \textbf{41} (1999), 236--266.

\bibitem{Connes}
A.~Connes,
``Trace formula in noncommutative geometry and the zeros of the Riemann
zeta function'',
\emph{Selecta Math.\ (N.S.)} \textbf{5} (1999), 29--106.

\bibitem{deBranges}
L.~de~Branges,
``The Riemann hypothesis for Hilbert spaces of entire functions'',
\emph{Bull.\ Amer.\ Math.\ Soc.} \textbf{15} (1986), 1--17.

\bibitem{Goldfeld}
D.~Goldfeld,
``A spectral interpretation of Weil's explicit formula'',
\emph{Lecture Notes in Math.} \textbf{1593} (1994), 135--152.

\bibitem{Lagarias}
J.C.~Lagarias,
``On a positivity property of the Riemann $\xi$-function'',
\emph{Acta Arithmetica} \textbf{89} (1999), 217--234.

\bibitem{Paper1}
U.~Tehrani,
``A Curvature Decomposition of the Explicit Formula'',
Zenodo, \url{https://doi.org/10.5281/zenodo.19025598}, 2026.

\bibitem{Paper2}
U.~Tehrani,
``From Local Curvature to the Weil Functional'',
Zenodo, \url{https://doi.org/10.5281/zenodo.19106992}, 2026.

\bibitem{Code}
U.~Tehrani,
Verification scripts for the prime--zero coupling series,
GitHub, \url{https://github.com/utehrani/analysislab-nt}.

\bibitem{Tate}
J.T.~Tate,
``Fourier analysis in number fields, and Hecke's zeta-functions'',
in: \emph{Algebraic Number Theory} (Cassels--Fr\"ohlich),
Academic Press, 1967.

\bibitem{Weil}
A.~Weil,
``Sur les `formules explicites' de la th\'eorie des nombres premiers'',
\emph{Comm.\ S\'em.\ Math.\ Univ.\ Lund, Suppl.}\ (1952), 252--265.

\end{thebibliography}


% --------------------------------------------------------
% Footer — normative format (v1.6)
% --------------------------------------------------------
\enlargethispage{3\baselineskip}
\vspace*{\fill}
\begin{minipage}{\linewidth}
\rule{\linewidth}{0.4pt}\smallskip\\
\small\emph{Paper~3 v3.2 \textperiodcentered\ May~2026
\textperiodcentered\ Pauca sed matura}
\end{minipage}

\end{document}
